%! Author = Saeid Yazdani
%! Date = 9/1/2024

\section{Code Listings of Important VHDL Modules}
\label{sec:hvs-fpga-compensation-vhdl}

\begin{lstlisting}[language=VHDL, caption={Compensation Control VHDL Module},
    label={lst:fpga-compensation-entity}]
entity compensation_ctrl is
port (
    clk                 : in  std_logic;
    rst_n               : in  std_logic;
    trigger             : in  std_logic;
    pre_comp_array      : in  t_comp_array  (0 to c_comp_width-1 );
    pa_comp_array       : in  t_comp_array  (0 to c_comp_width-1 );
    timing_array        : in  t_timing_array(0 to c_comp_width-1 );
    compensation_data   : OUT std_logic_vector(15 downto 0);
    PACmp_C             : OUT std_logic_vector(9 -1 downto 0);
    PACmp_R             : OUT std_logic_vector(7 -1 downto 0);
    PreCmp_C            : OUT std_logic_vector(8 -1 downto 0);
    PreCmp_R            : OUT std_logic_vector(8 -1 downto 0)
);
end entity;

architecture rtl of compensation_ctrl is
    signal s_pa_cmp_cr         : std_logic_vector(15 downto 0);
    signal s_pre_cmp_cr        : std_logic_vector(15 downto 0);
    signal s_trigger           : std_logic;

begin

    p_main : process (clk,rst_n)

        type states is (  idle, run, set, inc_vi, check_vi );
        variable state : states;
        variable v_cnt   : unsigned(c_tim_width-1 downto 0);
        variable pre_cnt : unsigned(4 downto 0);
        subtype t_pos is natural range 0 to c_comp_width;
        variable v_i : t_pos;

    begin

        if rst_n = '0' then

            s_pre_cmp_cr <= ResetValue_PREAMP_7 ;
            s_pa_cmp_cr  <= ResetValue_PAAMP_7  ;
            state        := idle;
            v_cnt        := (others => '0');
            pre_cnt      := (others => '0');
            v_i          := 0;
            s_trigger    <= '0';

        elsif rising_edge(clk) then

            s_trigger  <= trigger;

            case state is

                when idle =>

                    v_cnt        := (others => '0');
                    v_i          := 0;
                    s_pre_cmp_cr <=  pre_comp_array(c_comp_width-1);
                    s_pa_cmp_cr  <=  pa_comp_array (c_comp_width-1);

                    if  s_trigger = '0' and trigger = '1' then
                        state        := run;
                    end if;

                when run =>

                    if  s_trigger = '0' and trigger = '1' then

                        v_cnt        := (others => '0');
                        pre_cnt      := (others => '0');
                        v_i          := 0;
                        s_pre_cmp_cr <=  pre_comp_array(c_comp_width-1);
                        s_pa_cmp_cr  <=  pa_comp_array (c_comp_width-1);

                    else

                        if v_cnt >= unsigned(timing_array(v_i)) then

                            state := set;

                            s_pre_cmp_cr <= pre_comp_array(v_i);
                            s_pa_cmp_cr  <= pa_comp_array (v_i);

                        else
                            if pre_cnt = 19 then
                                v_cnt := v_cnt + 1;
                                pre_cnt := (others => '0');
                            else
                                pre_cnt := pre_cnt + 1;
                            end if;
                        end if;

                end if;

                when set =>

                        s_pre_cmp_cr <= pre_comp_array(v_i);
                        s_pa_cmp_cr  <= pa_comp_array (v_i);

                        state := inc_vi;

                when inc_vi =>

                    v_i := v_i +1;
                    state := check_vi;

                when check_vi =>

                    if v_i = c_comp_width then
                        state := idle;
                    else
                        state := run;
                    end if;

            end case;

        end if;

    end process;


    PreCmp_C   <= s_pre_cmp_cr(15 downto 8) ;
    PreCmp_R   <= s_pre_cmp_cr(7  downto 0) ;

    PACmp_C    <= s_pa_cmp_cr (15 downto 7) ;
    PACmp_R    <= s_pa_cmp_cr (6  downto 0) ;

    compensation_data <= s_pa_cmp_cr;

end architecture;
\end{lstlisting}

\newpage

\begin{lstlisting}[language=VHDL, basicstyle=\ttfamily\tiny, lineskip=-2pt,
    caption={ADC Calibration Module},
    label={lst:fpga-calibration-entity}]
entity adc_calibration is
    port (
        clk                 : in  std_logic;
        rst_n               : in  std_logic;
        -- adaptation
        offset              : in  std_logic_vector(15 downto 0);
        -- adc interface
        data_in             : in  std_logic_vector(15 downto 0);
        data_out            : out std_logic_vector(15 downto 0);
        data_en             : out std_logic
    );
end entity;

architecture rtl of adc_calibration is

    signal s_data_en   : std_logic;
    signal s_data      : signed(15 downto 0);
    signal s_result_s    : signed(16 downto 0);
    --signal s_result_u    : unsigned(17 downto 0);
    signal s_data_in   : std_logic_vector(15 downto 0);

begin

    s_data_in <= data_in;

    p_main : process (clk,rst_n)

        type states is ( idle, overflow );
        variable state : states;

    begin

        if rst_n = '0' then

            s_data_en  <= '0';
            s_data     <= (others => '0');
            s_result_s   <= (others => '0');
            --s_result_u   <= (others => '0');
            state := idle;

        elsif rising_edge(clk) then

            case state is

                when idle =>

                    s_result_s <=   signed( s_data_in(15) & s_data_in) +   signed( offset(15) & offset);
                    state := overflow;
                    s_data_en <= '0';

                when overflow =>
                    if s_result_s(16 downto 15) = "10"  then
                        s_data <= x"8000";
                    elsif s_result_s(16 downto 15) = "01" then
                        s_data <= x"7fff";
                    else
                        s_data <= s_result_s(15 downto 0);
                    end if;
                    state := idle;
                    s_data_en <= '1';

            end case;

        end if;

    end process;

    data_en    <= s_data_en;
    data_out   <= std_logic_vector(s_data);

end architecture;
\end{lstlisting}


\clearpage
\begin{lstlisting}[language=VHDL, basicstyle=\ttfamily\tiny, lineskip=-2pt,
    caption={Scaling Module},
    label={lst:fpga-scaling-entity}]
entity scaling is
    port (
        clk                 : in  std_logic;
        rst_n               : in  std_logic;
        -- adaptation
        scale               : in  std_logic_vector(15 downto 0);
        -- input interface
        data_in             : in  std_logic_vector(15 downto 0);
        conv_en             : in  std_logic;
        -- output interface
        data_out            : out std_logic_vector(15 downto 0);
        data_en             : out std_logic
    );
end entity;

architecture rtl of scaling is

    signal s_dac_en         : std_logic;
    signal s_dac_data       : signed(15 downto 0);
    signal s_result         : signed(33 downto 0);

begin

    p_main : process (clk,rst_n)

        type states is ( idle, overflow );
        variable state : states;

    begin

        if rst_n = '0' then

            s_result   <= (others => '0');
            state      := idle;
            s_dac_en   <= '0';
            s_dac_data <= (others => '0');

        elsif rising_edge(clk) then

            case state is

                when idle =>
		    s_dac_en   <= '0';
                    if conv_en = '1' then
                        s_result <= signed(data_in(15)&data_in) * signed('0'&scale);
                        state := overflow;
                    end if;

                when overflow =>
                    if s_result(33) = '1' and s_result(32 downto 27) /= "111111" then
                        s_dac_data <= "1000000000000000";
                    elsif s_result(33) = '0' and s_result(32 downto 27) /= "000000" then
                        s_dac_data <= "0111111111111111";
                    else
                        s_dac_data <= s_result(27 downto 12);
                    end if;
                    state := idle;
                    s_dac_en <= '1';

            end case;

        end if;

    end process;

    data_en     <= s_dac_en;
    data_out   <= std_logic_vector(s_dac_data);


end architecture;

\end{lstlisting}


\clearpage
\begin{lstlisting}[language=VHDL, basicstyle=\ttfamily\scriptsize, lineskip=-1pt,
    caption={Output Prepration (shared between COARSE and VERNIER)},
    label={lst:fpga-outputprep-entity}]

entity output_preparation is
generic ( DATA_WIDTH : natural := 16 );
port (
clk                 : in  std_logic;
rst_n               : in  std_logic;
-- interface
data_in             : in  std_logic_vector(DATA_WIDTH -1 downto 0);
scale               : in  std_logic_vector(DATA_WIDTH -1 downto 0);
offset              : in  std_logic_vector(DATA_WIDTH -1 downto 0);
data_out            : out std_logic_vector(DATA_WIDTH -1 downto 0);
data_en             : out std_logic
);
end entity;

architecture rtl of output_preparation is

signal s_data_sc2ca       : std_logic_vector(DATA_WIDTH -1 downto 0);

begin

inst_scale : entity rtl_lib.scaling
port map (
clk                 =>  clk         ,
rst_n               =>  rst_n       ,
-- adaptation
scale               =>  scale       ,
-- adc interface
data_in             =>  data_in     ,
conv_en             =>  '1'         ,
-- adc - dac interface
data_out            =>  s_data_sc2ca,
data_en             =>  open
);


inst_calib : entity rtl_lib.adc_calibration
port map (
clk                 => clk         ,
rst_n               => rst_n       ,
-- adaptation
offset              => offset      ,
-- dac interface    =>
data_in             => s_data_sc2ca,
data_out            => data_out    ,
data_en             => data_en
);

end architecture;
\end{lstlisting}




